\section{Predikat dan Variabel}

\logic{Logika predikat} memperluas logika proposisional dengan predikat dan kuantifier \cite{rosen2019,forallx}. Dalam logika proposisional, kita hanya dapat mengekspresikan pernyataan atomik sebagai kesatuan tanpa struktur internal; kita tidak dapat menangkap hubungan seperti ``semua'', ``ada'', atau ``untuk setiap''. Pernyataan seperti ``Semua manusia fana'' atau ``Ada bilangan prima genap'' tidak dapat direpresentasikan dengan baik dalam logika proposisional tanpa kehilangan struktur logisnya \cite{wiki_fol,copi2014}. Logika predikat mengatasi keterbatasan ini dengan memperkenalkan variabel dan predikat yang dapat mengambil objek dalam suatu domain sebagai argumen.

Gambar~\ref{fig:proposisional-vs-predikat} mengilustrasikan perbedaan konseptual antara logika proposisional dan logika predikat \cite{berkeley_fol,olp_fol}.

\begin{figure}[htbp]
\centering
\begin{tikzpicture}[node distance=1.2cm, font=\small]
  \node[proposisi, align=center] (prop) {\parbox{2cm}{\centering Proposisi atom $p$, $q$, $r$}};
  \node[proposisi, align=center, right=3.2cm of prop] (pred) {\parbox{2.2cm}{\centering Predikat + variabel $P(x)$, $Q(x,y)$}};
  \node[below=0.6cm of prop, align=center] (d1) {Tanpa struktur internal};
  \node[below=0.6cm of pred, align=center] (d2) {Domain + argumen};
  \node[operator, align=center, below=1.8cm of prop] (op1) {\parbox{2.35cm}{\centering Logika\\Proposisional}};
  \node[operator, align=center, below=1.8cm of pred] (op2) {\parbox{2.2cm}{\centering Logika\\Predikat}};
  \draw[arrow] (prop) -- (op1);
  \draw[arrow] (pred) -- (op2);
\end{tikzpicture}
\caption{Perbandingan konseptual: logika proposisional vs logika predikat}
\label{fig:proposisional-vs-predikat}
\end{figure}

Predikat adalah pernyataan yang mengandung satu atau lebih variabel; nilai kebenarannya bergantung pada nilai yang diberikan kepada variabel tersebut \cite{rosen2019,forallx}. Misalnya, ``$x$ lebih besar dari 5'' dapat ditulis sebagai $P(x)$ dengan $P(x)$ menyatakan ``$x > 5$''; maka $P(7)$ benar sedangkan $P(3)$ salah. Predikat memungkinkan kita merepresentasikan sifat atau relasi yang berlaku untuk sejumlah objek tanpa harus mendaftarkan semuanya secara eksplisit \cite{berkeley_fol,olp_fol}.

Predikat memiliki \emph{arity}, yaitu jumlah argumen yang diterima \cite{setslogic,copi2014}. Tabel~\ref{tab:arity-predikat} merangkum jenis predikat berdasarkan arity beserta contoh notasi dan makna.

\begin{table}[htbp]
\centering
\small
\begin{tabular}{|>{\raggedright\arraybackslash}p{1.15cm}|l|>{\raggedright\arraybackslash}p{3.2cm}|>{\raggedright\arraybackslash}p{2cm}|>{\raggedright\arraybackslash}p{3.5cm}|}
\hline
\textbf{Arity} & \textbf{Notasi} & \textbf{Makna contoh} & \textbf{Domain contoh} & \textbf{Nilai benar/salah} \\
\hline
Unary & $P(x)$ & ``$x$ adalah mahasiswa'' & Manusia & $P(\text{Ali})$ B, $P(\text{guru})$ S \\
\hline
Biner & $Q(x,y)$ & ``$x$ ibu kota dari $y$'' & Negara, kota & $Q(\text{Jakarta},\allowbreak\text{Indonesia})$ B \\
\hline
$n$-ary & $R(x_1,\ldots,x_n)$ & Relasi $n$ objek & Sesuai konteks & Tergantung interpretasi \\
\hline
\end{tabular}
\caption{Arity predikat: unary, biner, dan $n$-ary dengan contoh}
\label{tab:arity-predikat}
\end{table}

Predikat unary seperti $P(x)$ menyatakan sifat tunggal; predikat biner seperti $Q(x,y)$ menyatakan relasi antara dua objek. Predikat dengan arity lebih tinggi digunakan untuk relasi yang melibatkan tiga objek atau lebih. Variabel dapat bersifat bebas (free) atau terikat (bound); variabel terikat berada di dalam lingkup kuantifier yang akan dibahas di section berikutnya \cite{iep_predicate,olp_fol}.

Domain atau \emph{universe of discourse} adalah himpunan objek yang menjadi acuan untuk variabel \cite{rosen2019,forallx}. Pemilihan domain sangat mempengaruhi nilai kebenaran: predikat $P(x) = \text{``}x\text{ adalah bilangan prima''}$ bernilai berbeda jika domainnya bilangan bulat positif versus bilangan real. Tabel~\ref{tab:domain-predikat} memberi contoh domain dan predikat yang sesuai \cite{copi2014,stanford_cs103}.

\begin{table}[htbp]
\centering
\small
\begin{tabular}{|>{\raggedright\arraybackslash}p{2.2cm}|>{\raggedright\arraybackslash}p{3.8cm}|l|}
\hline
\textbf{Domain} & \textbf{Contoh predikat} & \textbf{Contoh nilai} \\
\hline
Bilangan bulat positif & $P(x)$: $x$ prima; $Q(x)$: $x > 5$ & $P(7)$ B, $P(4)$ S \\
\hline
Mahasiswa & $M(x)$: $x$ lulus; $R(x)$: $x$ rajin & $M(\text{Ana})$ B/S \\
\hline
Kota, negara & $I(x,y)$: $x$ ibu kota $y$ & $I(\text{Jakarta},\text{Indonesia})$ B \\
\hline
\end{tabular}
\caption{Contoh domain dan predikat}
\label{tab:domain-predikat}
\end{table}

Dalam matematika, domain sering kali bilangan bulat, bilangan real, atau himpunan tertentu; dalam konteks sehari-hari, domain bisa berupa mahasiswa, kota, atau entitas lain yang relevan \cite{copi2014,stanford_cs103}.

Dalam informatika, logika predikat memiliki penerapan langsung \cite{rosen2019,stanford_cs103}. Basis data relasional memodelkan data sebagai relasi; setiap relasi dapat dipandang sebagai predikat yang benar untuk tuple yang ada dalam tabel. Query SQL dengan klausa \code{WHERE} pada dasarnya mengekspresikan kondisi logika predikat. Representasi pengetahuan dalam sistem kecerdasan buatan sering menggunakan logika predikat orde satu (first-order logic) untuk menangani objek, relasi antar objek, dan kuantifikasi \cite{berkeley_fol,mit_6042}. Pemahaman predikat dan variabel menjadi fondasi untuk memahami query basis data, reasoning engine, dan spesifikasi formal perangkat lunak.

\begin{contoh}
\textbf{Evaluasi predikat dengan Python} \cite{python_docs_boolean,rosen2019}

Fungsi berikut mengevaluasi predikat $P(x) = \text{``}x > 5\text{''}$ pada domain diskret (daftar bilangan) dan mencetak pasangan $(x, P(x))$. Ide \code{all(...)} dan \code{any(...)} nantinya bersesuaian dengan kuantifier $\forall$ dan $\exists$.

\begin{lstlisting}[style=pythonstyle, caption={Evaluasi predikat pada domain diskret}]
def eval_predicate(domain, predicate):
    """Cetak (elemen, nilai predikat) untuk setiap elemen domain."""
    for x in domain:
        print(f"  P({x}) = {predicate(x)}")

# Predikat: x > 5
P = lambda x: x > 5
domain = [3, 5, 7, 10]
eval_predicate(domain, P)   # P(3)=False, P(5)=False, P(7)=True, P(10)=True

# Kuantifier (akan dibahas di section berikutnya):
# forall: all(P(x) for x in domain)  -> False
# exists: any(P(x) for x in domain)  -> True
\end{lstlisting}
\end{contoh}

\begin{contoh}
``$x$ lebih besar dari 5'' ditulis $P(x)$ dengan $P(x) = \text{``}x > 5\text{''}$. Maka $P(7)$ benar dan $P(3)$ salah. Untuk predikat biner ``$x$ adalah ibu kota dari $y$'', misal $Q(x,y)$ dengan domain negara dan kota, maka $Q(\text{Jakarta}, \text{Indonesia})$ benar dan $Q(\text{Bandung}, \text{Indonesia})$ salah.
\end{contoh}

\begin{konsep}
Predikat adalah pernyataan dengan variabel yang nilainya bergantung pada argumen. Domain menentukan himpunan objek yang diacu. Predikat unary menyatakan sifat; predikat biner atau lebih menyatakan relasi. Variabel bebas belum terikat kuantifier; variabel terikat berada dalam lingkup $\forall$ atau $\exists$.
\end{konsep}
